% БҮЛЭГ 2: ХОЛБОГДОХ СУДАЛГАА БА ОНОЛЫН ҮНДЭС
% 8-бүлгийн шинэ бүтцийн дагуу: онолын үндэс + related works + similar systems
%
\chapter{Холбогдох судалгаа ба онолын үндэс}
\label{Chapter2}

Энэхүү бүлэгт NLP, мэдрэмжийн шинжилгээ, хортой агуулга илрүүлэлт,
Transformer/MN-BERT онолын үндсийг тайлбарлаж, холбогдох судалгааны орхигдсон
чиглэлийг тодорхойлно.

%===============================================================================
%  ОНОЛЫН ХЭСЭГ
%===============================================================================

%-------------------------------------------------------------------------------
\section{Эх хэлний боловсруулалт (NLP)}

NLP нь хүний хэлийг компьютерт ойлгуулах, тайлбарлах хиймэл оюун ухааны дэд салбар бөгөөд токенчлол, нормчлол, векторжуулалт (TF-IDF, Word Embedding) зэрэг үе шатуудтай. Монгол хэл нь залгамал бүтэцтэй тул Transformer-д суурилсан загварууд илүү тохиромжтой (\ref{sec:transformer}-р хэсэгт).

%-------------------------------------------------------------------------------
\section{Мэдрэмжийн шинжилгээ (Sentiment Analysis)}

Мэдрэмжийн шинжилгээ нь текстийн хандлагыг автоматаар тодорхойлох NLP-ийн дэд чиглэл \cite{liu2012}. Дүрмэд суурилсан, машин сургалтад суурилсан (NB/SVM/LR + TF-IDF), гүн сургалтад суурилсан (Transformer) гэсэн гурван арга зүйтэй; энэ судалгаанд өгүүлбэр/баримт бичгийн түвшний (sentence/document-level) шинжилгээ ашиглана.

%-------------------------------------------------------------------------------
\section{Хортой агуулга илрүүлэх систем}

Хортой агуулга илрүүлэх нь онлайн доромжлол, дарамт, заналхийлэлийг автоматаар таних NLP-ийн тусгайлсан салбар (Jigsaw Perspective API 2016, Kaggle 2018). Гол бэрхшээл: контекстийн нарийн байдал, егөөдөл/ёжлол; Монгол хэлийг корпусын хомсдол, тэмдэглэгээний (annotation) стандартчилал дутуу байдал нэмэгдүүлнэ. Монгол корпус дутмаг тул гар аргаар шошголох шаардлагатай.

%-------------------------------------------------------------------------------
\section{Бүтээлч шүүмжлэл ба ялгах шалгуур}

Бүтээлч шүүмжлэл нь гадаргуугаас сөрөг мэт харагддаг ч
аргументтай, сайжруулах санал агуулсан, хувь хүн рүү биш асуудал руу чиглэсэн
хандлагын тусгай хэлбэр юм. Автоматаар ялгах зорилгоор хоёрын хоорондын
шалгуурыг тодорхой тогтоох шаардлагатай. Доорх хүснэгт нь хоёр ангиллын
үндсэн ялгааг харуулна.

\begin{table}[ht]
  \centering
  \small
  \caption{Бүтээлч шүүмжлэл ба хортой сөрөг агуулгын харьцуулалт}
  \label{tab:constructive-vs-toxic}
  \begin{tabular}{@{}lp{4.6cm}p{4.6cm}@{}}
    \toprule
    \textbf{Шалгуур} & \textbf{Бүтээлч шүүмжлэл} & \textbf{Хортой сөрөг агуулга} \\
    \midrule
    Чиглэл        & Асуудал, үзэгдэл рүү чиглэсэн      & Хувь хүн рүү чиглэсэн \\[2pt]
    Аргументаци   & Логик, нотолгоо агуулсан            & Сэтгэл хөдлөлд суурилсан, нотолгоогүй \\[2pt]
    Санал агуулга & Сайжруулах арга замыг тусгасан      & Шийдэл санал болгоогүй \\[2pt]
    Хэлний хэлбэр & Сөрөг ч таарсан, агуулга хэрэгтэй  & Доромжлол, заналхийлэл \\[2pt]
    Зорилго       & Нөхцөл байдлыг сайжруулах           & Гомдоох, дарамтлах, гутаах \\
    \bottomrule
  \end{tabular}
\end{table}

Ялган ангилах нь NLP-ийн хамгийн бэрхшээлтэй даалгавруудын нэг: хоёулаа
сөрөг лексик агуулдаг тул лексик шинжилгээ хангалтгүй --- контекст, дискурс
гүнзгий ойлгох шаардлагатай тул Transformer загварууд хамгийн тохиромжтой.

%-------------------------------------------------------------------------------
\section{Нийгэм-сэтгэл зүйн онолын хүрээ: TAM загвар}

Хэрэглэгчийн хүлээн зөвшөөрлийн үнэлгээнд \textbf{TAM} \cite{davis1989}
(\textit{Perceived Usefulness} + \textit{Perceived Ease of Use})-ийг ашиглав.
``Бүтээлч шүүмжлэл'' ангиллыг тусад нь ялгах нь хэт автомат цензураас
зайлсхийх замаар нийгмийн хүлээн зөвшөөрлийг дэмжинэ (бүлэг~\ref{Chapter8}).

%-------------------------------------------------------------------------------
\section{Transformer архитектур ба BERT загвар}
\label{sec:transformer}

Transformer архитектур \cite{vaswani2017} нь Self-Attention механизмаар бүх
токен хоорондын хамаарлыг нэгэн зэрэг тооцоолж (RNN-ийг давж), Multi-Head
Attention нь үүнийг параллель толгойд давтдаг. Self-Attention механизм нь өгүүлбэрийн токен бүр бусад бүх токентой хэрхэн хамаарч байгааг жинлэж (weighting), контекст мэдээллийг гүнзгий тусгадаг. BERT \cite{devlin2019} нь хоёр чиглэлт (Bidirectional) Masked LM (өгүүлбэр дэх зарим үгийг халхалж таах) болон NSP (дараагийн өгүүлбэрийг таах) pre-training-аар их корпус дээр сургагдаж, fine-tuning-ээр NLP benchmark-ийг мэдэгдэхүйц сайжруулсан.

%-------------------------------------------------------------------------------
\subsection{Үг векторжуулалт (Word Embedding)}
\label{sec:word-embedding}

Машин сургалтын загварт текстийг тоон хэмжигдэхүүнд хөрвүүлэх шаардлагатай. BoW/TF-IDF нь контекстийн мэдрэмжгүй бөгөөд Монгол хэлний залгамал
бүтцэд үгийн сан (vocabulary) огцом томорч, нягтаршил буурдаг. Word2Vec/GloVe/FastText зэрэг \textit{static embedding} аргууд нь семантик ойролцоог тусгадаг ч нэг үгэнд нэг л вектор харгалзуулдаг тул полисемийг (нэг үг олон утгатай байх) тооцдоггүй. BERT зэрэг \textit{contextual embedding} нь векторыг нөхцөл байдлаас хамааруулан динамикаар тооцдог тул Монгол хэлний нийлмэл морфологийн семантик ялгааг илүү нарийн барина.
%-------------------------------------------------------------------------------
\section{MN-BERT --- Монгол хэлний загвар}

MN-BERT нь BERT-base (12 layer, 768 hidden, 12 head) архитектурт суурилсан,
Монгол Wikipedia-д pre-train хийгдсэн, HuggingFace-д нээлттэй загвар бөгөөд
SentencePiece токенчлолоор (tokenization) залгамал морфологийг задалдаг --- Монгол
NLP-ийн суурь загвар. Нарийн тохируулга явцад ангиллын толгой (classification head) нэмж, тодорхой даалгаварт --- манай тохиолдолд дөрвөн ангиллын систем --- дахин сургана.

MN-BERT-ийн нарийн тохируулгын архитектурыг зураг~\ref{fig:mnbert-arch}-д харуулав. Оролтын текстийг токенчилсний дараа \texttt{[CLS]} токений гаралтын вектор ангиллын толгойд дамжуулж, дөрвөн ангиллын нэгт олон ангиллын шийдвэр гаргадаг.

\begin{figure}[H]
\centering
\begin{tikzpicture}[
  box/.style={draw, rounded corners, minimum width=5.5cm, minimum height=0.75cm,
              align=center, font=\small},
  arrow/.style={->, >=Stealth, thick}
]
\node[box, fill=gray!15]  (input) at (0,0)
  {Оролтын текст (кирилл токенууд)};
\node[box, fill=yellow!25] (embed) at (0,-1.4)
  {Token + Position + Segment Embedding};
\node[box, fill=blue!15, align=center, minimum height=1.9cm] (trans) at (0,-3.15)
  {$12\times$ Transformer Encoder\\(Multi-Head Attention + Feed-Forward)};
\node[box, fill=green!20]  (cls)   at (0,-5.1)
  {\texttt{[CLS]} гаралтын вектор ($\mathbb{R}^{768}$)};
\node[box, fill=orange!25] (head)  at (0,-6.4)
  {Ангиллын толгой (Linear + Softmax)};
\node[box, fill=red!15, minimum width=7cm] (out) at (0,-7.65)
  {Эерэг \enspace Саармаг \enspace Хортой сөрөг \enspace Бүтээлч шүүмжлэл};
\draw[arrow] (input) -- (embed);
\draw[arrow] (embed) -- (trans);
\draw[arrow] (trans) -- (cls);
\draw[arrow] (cls)   -- (head);
\draw[arrow] (head)  -- (out);
\end{tikzpicture}
\caption{MN-BERT загварын архитектур (нарийн тохируулгын дараах)}
\label{fig:mnbert-arch}
\end{figure}

Ч.~Жаргалмаа (ШУТИС, 2025) \textit{``Нийгмийн сүлжээний платформуудын
...
хэрэглэгчийн зан төлөв хандлага болон тренд анализ''} бакалаврын судалгааны
ажилдаа~\cite{jargalmaa2025} MN-BERT нарийн тохируулгатай загвар болон уламжлалт
машин сургалтын суурь загваруудын гүйцэтгэлийн харьцуулалтыг тусгасан бөгөөд
тус үр дүнгийг доорх хүснэгтэд толилуулав.

\begin{table}[ht]
  \centering
  \caption{Загваруудын гүйцэтгэлийн харьцуулалт~\cite{jargalmaa2025}}
  \label{tab:model-comparison}
  \begin{tabular}{lcccc}
    \toprule
    \textbf{Загвар} & \textbf{Нарийвчлал} & \textbf{$F_1$ (жигн.)} &
    \textbf{Precision} & \textbf{Recall} \\
    \midrule
    MN-BERT + Нарийн тохируулга        & $82.7\%$ & $0.81$ & $0.83$ & $0.81$ \\
    Logistic Regression + TF-IDF       & $74.0\%$ & $0.71$ & $0.78$ & $0.79$ \\
    Naive Bayes + TF-IDF               & $70.0\%$ & $0.68$ & $0.75$ & $0.84$ \\
    SVM (LinearSVC) + TF-IDF           & $73.5\%$ & $0.72$ & $0.80$ & $0.82$ \\
    \bottomrule
  \end{tabular}
  \par\vspace{4pt}
  \begin{minipage}{\linewidth}
    \footnotesize\textit{Эх сурвалж:} Ч.~Жаргалмаа (2025).
    \textit{Нийгмийн сүлжээний платформуудын хэрэглэгчийн зан төлөв хандлага
    болон тренд анализ.} Бакалаврын дипломын ажил,
    Шинжлэх Ухаан, Технологийн Их Сургууль (ШУТИС)~\cite{jargalmaa2025}.
  \end{minipage}
\end{table}

\medskip
Хүснэгт~\ref{tab:model-comparison}-оос MN-BERT ($82.7\%$ нарийвчлал, $F_1{=}0.81$)
нь хамгийн ойрын LR+TF-IDF ($74.0\%$)-аас 8.7 нэгжээр давж байна ---
MN-BERT-ийг дөрвөн ангилалт системд ашиглах шийдлийг зөвтгөж байна.

%===============================================================================
%  ХОЛБОГДОХ СУДАЛГААНЫ ХЭСЭГ
%===============================================================================

%-------------------------------------------------------------------------------
\section{Холбогдох судалгааны ажлын шинжилгээ}
\label{sec:related-work}

Олон улсад хортой агуулга ангилах чиглэлд эрчимтэй судалгаа хийгдсэн:
(1)~суурь судалгаа; (2)~BERT-д суурилсан загвар; (3)~бага хэрэглэгддэг хэлний NLP;
(4)~Монгол хэлний NLP. Энэхүү шинжилгээгээр авч үзээгүй талуудыг тодорхойлно.

%-------------------------------------------------------------------------------
\subsection{Сөрөг агуулга илрүүлэлтийн суурь судалгаа}

Davidson et al. (2017) Twitter-ийн 24,802 нийтлэлийг crowd-sourcing аргаар
гурван ангилалд шошголж (Hate Speech, Offensive, Normal), LR/SVM-ээр
$90.8\%$ нарийвчлал гаргасан \cite{davidson2017} --- ``сөрөг'' агуулга дотоодоо
ялгалт шаардлагатайг нотолсон нь манай хоёр шатлалт архитектурын онолын үндэс.

%-------------------------------------------------------------------------------
\subsection{BERT болон Transformer-д суурилсан загваруудын сошиал медиа хортой агуулга ангилалт}

Caselli et al. (2021)-ийн HateBERT нь домайнд тохирсон pre-training-ийн ач
холбогдлыг нотолсон \cite{caselli2021} --- MN-BERT сонголтын онолын баталгаа.
Mathew et al. (2021)-ийн HateXplain-д олон шошгологч + Cohen's Kappa нийцэл
ашигласан нь манай аннотацийн загвар \cite{mathew2021};
Fan et al. (2021) BERT-base AUC-ROC 0.9856 гарсныг нотолсон \cite{fan2021toxic}.

%-------------------------------------------------------------------------------
\subsection{Бага хэрэглэгддэг хэлний NLP болон хөндлөн хэлний дайрах хэл
таних судалгаа}

Ranasinghe \&\ Zampieri (2021) долоон бага хэрэглэгддэг хэлд monolingual загвар
давуу болохыг нотолсон \cite{ranasinghe2021} --- MN-BERT сонголтын үндэс.
Nabiilah et al. (2023)-ийн IndoBERT нь mBERT-ийг $F_1$-ээр 17 нэгжээр давсан
\cite{nabiilah2023bert}; Akhter et al. (2022) Урду хэлд mBERT-ийн
давуу тал ($F_1{=}0.8149$)-ыг нотолсон \cite{akhter2022multi}.

%-------------------------------------------------------------------------------
\subsection{Авч үзээгүй талууд ба энэхүү дипломын ажлын оруулах хувь нэмэр}

Дээрх системчилсэн шинжилгээнээс таван орхигдсон чиглэл ажиглагдав:
(1)~Монгол NLP-д Хортой--Бүтээлч ялгалт байхгүй;
(2)~MN-BERT 4-ангилалт жишиг сорил (benchmark) нийтэд байхгүй;
(3)~хоёр шатлалт архитектур Монголд туршигдаагүй;
(4)~нийтэд нээлттэй 4-ангилалт корпус (corpus) байхгүй;
(5)~Монгол хэл бага хэрэглэгддэг NLP сорилуудад байхгүй.
Энэхүү ажил эдгээрийг нөхнө: (i)~4-ангилалт датасет бүрдүүлэх;
(ii)~хоёр шатлалт MN-BERT архитектур боловсруулах;
(iii)~залгамал онцлогт урьдчилсан боловсруулалт (preprocessing) хэрэгжүүлэх;
(iv)~уламжлалт машин сургалт болон 3-ангилалт MN-BERT-тэй харьцуулах;
(v)~датасетийг нийтэд нэвтрүүлэх.

%===============================================================================
%  ИЖИЛ ТӨСТЭЙ СИСТЕМИЙН СУДАЛГАА
%===============================================================================

%-------------------------------------------------------------------------------
\section{Ижил төстэй системийн судалгаа}

\subsection{Монгол хэлний мэдрэмжийн шинжилгээний ажлууд}

МУИС-ийн NLP багийн 2021 оны судалга \cite{mnunlp2021} нь 5000 текстийг
гурван ангилалд шошголж NB ($74\%$) ба SVM ($78\%$) загварыг туршсан суурь
ажил; гэвч гүн сургалтын загвар, 4 ангиллын ялгалт дутагдаж байна.
\cite{jargalmaa2025} MN-BERT-ээр $82.7\%$ нарийвчлал ($F_1{=}0.81$) гаргасан ч
гурван ангилалтад хязгаарлагдсан --- энэхүү ажил тэр ялгалтыг нөхнө.


\subsection{Олон улсын хортой агуулга илрүүлэх системүүд}

Perspective API (Jigsaw/Google, 2016) нь олон хэлний хортой агуулга илрүүлэх
арилжааны API бөгөөд нейрон сүлжээнд тулгуурладаг; Монгол зэрэг бага
хэрэглэгддэг хэлүүдэд дэмжлэг хязгаарлагдмал. Kaggle-ийн Toxic Comment
Challenge (2018) шилдэг шийдлүүд LSTM/CNN/Transformer хослол, multi-label
ангиллыг ашигласан нь манай архитектурт лавлагаа болно.

\subsection{Монгол хэлний платформуудтай харьцуулалт}

Sumbee.mn нь маркетинг/брэнд хяналтад, Brand24 болон Talkwalker нь Монгол хэлний дэмжлэгийг хэсэгчлэн санал болгодог. Гэвч эдгээр нь хортой агуулгыг бүтээлч шүүмжлэлээс ялгах тал дээр хязгаарлалттай.

Доорх хүснэгтэд санал болгож буй системийг одоо байгаа шийдлүүдтэй харьцуулав.

\begin{table}[H]
    \centering
    \caption{Одоо байгаа системүүд болон санал болгож буй системийн харьцуулалт}
    \label{tab:system-comparison}
    \begin{tabular}{@{}lccc@{}}
        \toprule
        \textbf{Үйлдэл / Шалгуур} & \textbf{Sumbee.mn} & \textbf{Brand24} & \textbf{Санал болгож буй} \\
        \midrule
        Polarity (Pos/Neg/Neu)       & $\checkmark$ & $\checkmark$ & $\checkmark$ \\
        Хортой сөрөг ялгалт          & $\times$     & $\times$     & $\checkmark$ \\
        Бүтээлч шүүмжлэл ялгалт     & $\times$     & $\times$     & $\checkmark$ \\
        Монгол хэлний морфологи      & Дунд         & Бага         & Өндөр (MN-BERT) \\
        Өсвөр насны аюулгүй байдал   & $\times$     & $\times$     & $\checkmark$ \\
        Идэвхтэй шүүлт (Block)       & $\times$     & $\times$     & $\checkmark$ \\
        \bottomrule
    \end{tabular}
    \par\vspace{4pt}
    \begin{minipage}{\linewidth}
      \footnotesize\textit{Эх сурвалж:} Зохиогчийн бэлтгэсэн харьцуулалт.
      Sumbee.mn болон Brand24-ийн мэдээлэл тухайн системийн нийтийн баримтаас авав.
    \end{minipage}
\end{table}
\medskip
Манай системийн давуу тал: ``Сөрөг'' ангиллыг нарийвчлан ялгаж, идэвхтэй шүүлт хийнэ (Хүснэгт~\ref{tab:system-comparison}).


%-------------------------------------------------------------------------------
\section{Бүлгийн дүгнэлт}

Transformer/BERT нь контекстийн ялгааг барих хамгийн дэвшилтэт хэрэгсэл;
MN-BERT Монгол NLP-д хамгийн үр ашигтай. Олон улсын ажлуудад хортой--бүтээлч
ялгалт нотлогдсон ч Монгол хэлээр тусгайлсан систем дутагдалтай ---
энэхүү ажил тэр орхигдсон чиглэлийг нөхнө.
